CS 学术写作规范指南
====================

参考来源：https://people.eecs.berkeley.edu/~pattrsn/talks/writingtips.html


一、段落结构 (Paragraph Structure)
----------------------------------

1.1 段落开头
- 首句承上启下：任何段落的第一句话需要同时与上一段的首句和尾句衔接（考虑两类读者：通读全文者和跳读者）
- 首句放核心信息：每个段落的第一句话要放最核心的信息

1.2 段落内部
- 句子过渡流畅：段落内，句子与句子之间要 transition well
- 信息不要交叉：不要 interleave information，相关信息要放在一起（不要 A→B→A→B，应该 A→A→B→B）


二、句子与信息单元 (Sentence & Information Unit)
------------------------------------------------

2.1 单一性原则
- 每个段落只能有一个核心信息
- 每个句子只能有一个核心信息

2.2 自解释原则 (Self-explained)
- 每个段落和每个句子要 self-explained
- 如果写了一个 statement，必须给出理由
- 支撑方式：facts / reasoning / facts + reasoning

2.3 表达要生动
- 对于 statement，可以用例子做补充解释
- 例子/事件要 cite paper，为 statement 提供更多 context


三、语法与措辞 (Grammar & Wording)
----------------------------------

3.1 语态与结构
- 少用被动语态
- 重要信息前置：句子中重要的信息放前面
- 句子要平衡
- 避免连续的多个形容词或名词短语并列
- 少用代词 it, they, this, that

3.2 术语使用
- 先定义后使用：在使用 terminology 之前要先定义
- 避免模糊用词：不要用 vague 和 unclear terminology
- 保持一致性：对同一个 point，用同样的词

3.3 主语变化
- 避免连续使用同一个主语，特别是 "we"
- 但重要的句子需要强调时，该用还是要用

3.4 数字规范
- 1-10 用拼写（one, two, three...）
- >11 用数字（12, 13, 14...）


四、文章结构 (Paper Structure)
------------------------------

4.1 内容定位
- 文章里任何一句话的内容和所在位置都需要有理由

4.2 章节设计
- 每章要和 Introduction match
- 如果 section 只有一段，不要单独给 section，用 textbold 代替
- 聚焦常见情况：Do not focus on corner case when writing paper, make the common case complete（避免遗漏 key spirit）

4.3 教育性
- Paper 需要 educational，让读者感觉学到了东西


五、Contribution 写作 (Writing Contributions)
---------------------------------------------

- 按重要性排序：contribution 按照重要性从高到低写
- 体现 novelty：contribution 要体现 novelty，而不是做了什么 labor work


六、Outline 写作规范 (Outline Writing)
--------------------------------------

- 分 section 写：每个 section 单独写 outline，在前一个 section 得到 feedback 之后再写下一个
- 完整且有序：outline 要包含所有点的信息，要按顺序写


总结速查表
==========

类别          关键点
----          ------
段落开头      首句承上启下 + 首句放核心
段落内部      过渡流畅 + 信息集中不交叉
信息单元      单一性 + 自解释 + 有支撑
表达技巧      用例子 + cite paper
语态结构      少被动 + 重要前置 + 句子平衡
术语使用      先定义 + 不模糊 + 保持一致
主语变化      避免重复"we"，但重要时例外
数字规范      1-10拼写，>11数字
文章结构      内容与位置有理由 + 章节match intro
Section       单段不设section，用bold
核心精神      聚焦common case，不纠结corner case
教育性        让读者有收获
Contribution  重要性排序 + 体现novelty
Outline       分section写 + 完整有序